% =====================================================================
%  Banco complementar --- Campo Eletrico  (REVISADO)
%  Substitui a versao gerada por template (8 clones em N1, 11 em N2,
%  7 em N3 e 8 questoes conceituais marcadas como N4).
%  O cap01b e denso e ja cobre: E de carga puntiforme, proporcionalidade
%  com a distancia, carga de prova com forca dada, direcao e sentido do
%  campo, PONTO DE CAMPO NULO (tres vezes: Q e 4Q; 9 e 4 uC; -Q e +4Q),
%  duas cargas de mesmo modulo a 20 cm, grafico E x d, ELETRON ENTRE
%  PLACAS PARALELAS com deflexao, quadrado de quatro cargas, anel
%  isolante no eixo, haste finita por integracao (N4) e momento dipolar
%  (N4).
%  Este banco evita esses casos e explora: E=0 com V nao nulo, linhas de
%  campo e por que nao se cruzam, PROJETO de placas paralelas, cavidade
%  em condutor, tracado numerico de linhas, esfera isolante macica por
%  Gauss, par de placas com +-sigma e o limite imposto pela rigidez
%  dieletrica.
%  Distribuicao mantida: 8 N1 + 11 N2 + 7 N3 + 8 N4 = 34
% =====================================================================

\subsubsection*{Banco complementar --- Campo elétrico}

\begin{atencao}[Organização do banco]
O campo é \textbf{vetorial}: contribuições somam-se por componentes, e podem
cancelar-se. Ele é definido por $\vec E=\vec F/q_0$, mas \emph{não depende} da
carga de prova --- é propriedade do ponto. Duas expressões dominam o tópico:
$E=\ke Q/r^{2}$ para carga puntiforme e $E=\sigma/\varepsilon_0$ para plano
carregado, esta última \emph{independente da distância}. O N4 trata de projeto,
condutores, distribuições contínuas e limites físicos.
\end{atencao}

\blocofixacao

\begin{questao}[1]{}
Determine o módulo do campo elétrico produzido por $Q=+\SI{3}{\micro\coulomb}$
a $\SI{30}{\centi\meter}$ de distância.
\dados{$\ke=\SI{9e9}{\newton\meter\squared\per\coulomb\squared}$}
\resposta{$E=\SI{300}{\kilo\newton\per\coulomb}$}
\resolucao{
\[
E=\frac{\ke Q}{r^{2}}=\frac{\pot{9}{9}\times\pot{3}{-6}}{(0{,}30)^{2}}
=\frac{\pot{2.7}{4}}{0{,}09}=\pot{3}{5}\,\si{\newton\per\coulomb}.
\]
O campo aponta \textbf{para longe} da carga, por ela ser positiva.\\
\textbf{Atenção à potência:} o campo cai com $1/r^{2}$, o potencial com
$1/r$. Dobrar a distância divide $E$ por quatro e $V$ por dois.}
\end{questao}

\begin{questao}[1]{}
Determine o campo produzido por $Q=-\SI{5}{\micro\coulomb}$ a
$\SI{50}{\centi\meter}$, indicando o sentido.
\resposta{$E=\SI{180}{\kilo\newton\per\coulomb}$, apontando \emph{para} a carga}
\resolucao{
\[
E=\frac{\pot{9}{9}\times\pot{5}{-6}}{0{,}25}
=\pot{1.8}{5}\,\si{\newton\per\coulomb}.
\]
\textbf{Sentido:} cargas negativas \emph{atraem} cargas de prova positivas, logo
o campo aponta \textbf{para} a carga fonte.\\
\textbf{Regra mnemônica:} o campo \emph{sai} do positivo e \emph{entra} no
negativo --- é a mesma regra que orienta as linhas de campo.}
\end{questao}

\begin{questao}[1]{}
Uma carga $q=+\SI{2}{\micro\coulomb}$ é colocada em um ponto onde
$E=\SI{5000}{\newton\per\coulomb}$. Determine a força sobre ela.
\resposta{$F=\SI{10}{\milli\newton}$, no sentido do campo}
\resolucao{
\[
F=qE=\pot{2}{-6}\times5000=\pot{1}{-2}\,\si{\newton}=\SI{10}{\milli\newton}.
\]
Como a carga é positiva, a força tem o \emph{mesmo} sentido de $\vec E$. Se
fosse negativa, teria sentido oposto --- e apenas o sentido mudaria, não o
módulo.}
\end{questao}

\begin{questao}[1]{}
Determine a carga que produz $E=\SI{2000}{\newton\per\coulomb}$ a
$\SI{30}{\centi\meter}$.
\resposta{$Q=\SI{20}{\nano\coulomb}$}
\resolucao{
\[
Q=\frac{Er^{2}}{\ke}=\frac{2000\times0{,}09}{\pot{9}{9}}
=\pot{2}{-8}\,\si{\coulomb}=\SI{20}{\nano\coulomb}.
\]
Uma carga de nanocoulombs --- ordem de grandeza típica da eletrização por
atrito --- já produz milhares de newtons por coulomb a trinta centímetros.}
\end{questao}

\begin{questao}[1]{}
A que distância de $Q=+\SI{8}{\micro\coulomb}$ o campo vale
$\SI{200}{\kilo\newton\per\coulomb}$?
\resposta{$r=\SI{0.6}{\meter}$}
\resolucao{
\[
r=\sqrt{\frac{\ke Q}{E}}=\sqrt{\frac{\pot{9}{9}\times\pot{8}{-6}}{\pot{2}{5}}}
=\sqrt{0{,}36}=\SI{0.6}{\meter}.
\]
\textbf{Superfície de campo constante:} todos os pontos a
$\SI{0.6}{\meter}$ da carga têm o mesmo \emph{módulo} de campo --- mas
\emph{direções} diferentes, pois o campo é radial. Diferentemente das
equipotenciais, essas superfícies não têm significado especial.}
\end{questao}

\begin{questao}[1]{}
O campo elétrico em um ponto, medido com carga de prova $q_0$, vale $E$. Se a
carga de prova for \textbf{dobrada}, o valor medido:
\begin{alternativas}[2]
\item dobra
\item cai à metade
\item não se altera
\item depende do sinal de $q_0$
\end{alternativas}
\resposta{(c)}
\resolucao{Por definição, $\vec E=\vec F/q_0$. Dobrar $q_0$ dobra a força
\emph{e} o denominador --- a razão permanece.\\
\textbf{Ponto conceitual:} o campo é uma propriedade \textbf{do ponto do
espaço}, criada pelas cargas-fonte. A carga de prova apenas o \emph{revela};
ela não o cria nem o modifica.\\
\textbf{Ressalva real:} se $q_0$ for grande demais, ela perturba as
cargas-fonte (por indução, em condutores) e a medida deixa de valer. Daí a
definição rigorosa usar o limite $q_0\to0$.}
\end{questao}

\begin{questao}[1]{}
Mostre que $\si{\newton\per\coulomb}$ e $\si{\volt\per\meter}$ são a mesma
unidade.
\resposta{$\si{\volt}=\si{\joule\per\coulomb}$ e
$\si{\joule}=\si{\newton\meter}$}
\resolucao{
\[
\frac{\si{\volt}}{\si{\meter}}
=\frac{\si{\joule}}{\si{\coulomb}\cdot\si{\meter}}
=\frac{\si{\newton}\cdot\si{\meter}}{\si{\coulomb}\cdot\si{\meter}}
=\frac{\si{\newton}}{\si{\coulomb}}.\ \checkmark
\]
\textbf{As duas leituras do campo.} $\si{\newton\per\coulomb}$ enfatiza a
\emph{força} por unidade de carga; $\si{\volt\per\meter}$ enfatiza a
\emph{taxa de queda do potencial} por unidade de comprimento. São descrições
equivalentes de $\vec E=-\nabla V$.\\
Na prática, usa-se $\si{\volt\per\meter}$ quando há tensões envolvidas (placas,
cabos) e $\si{\newton\per\coulomb}$ em problemas de dinâmica de partículas.}
\end{questao}

\begin{questao}[1]{}
Uma carga negativa é colocada em uma região de campo elétrico uniforme. Sobre
seu movimento inicial, é correto afirmar que ela:
\begin{alternativas}[2]
\item acelera no sentido do campo
\item acelera no sentido oposto ao campo
\item permanece em repouso
\item move-se perpendicularmente ao campo
\end{alternativas}
\resposta{(b)}
\resolucao{A força vale $\vec F=q\vec E$ com $q<0$, logo $\vec F$ é
\textbf{antiparalela} a $\vec E$.\\
\textbf{Consequência energética:} a carga negativa move-se espontaneamente no
sentido de $V$ \emph{crescente} --- o oposto de uma carga positiva. Em ambos os
casos, a energia potencial $U=qV$ \emph{diminui}, como deve ocorrer em qualquer
movimento espontâneo.}
\end{questao}

\blocoaplicacao

\begin{questao}[2]{}
Uma carga $q_0=-\SI{1}{\micro\coulomb}$ é colocada em um campo uniforme de
$E=\SI{2000}{\newton\per\coulomb}$.
\begin{itensq}
\item Determine a força.
\item Determine sua aceleração, se a massa for $\SI{4}{\gram}$.
\end{itensq}
\resposta{$F=\SI{2}{\milli\newton}$, oposta a $\vec E$;
$a=\SI{0.5}{\meter\per\second\squared}$}
\resolucao{(a) $F=|q|E=\pot{1}{-6}(2000)=\pot{2}{-3}\,\si{\newton}$, com sentido
\emph{oposto} ao campo.\\
(b)
\[
a=\frac{F}{m}=\frac{\pot{2}{-3}}{\pot{4}{-3}}=\SI{0.5}{\meter\per\second\squared}.
\]
\textbf{Comparação com a gravidade:} $a$ é apenas $5\%$ de $g$. Para objetos
macroscópicos, a força elétrica só compete com o peso se a carga for
relativamente grande ou a massa muito pequena --- daí os experimentos
eletrostáticos clássicos usarem pedaços de papel e gotas de óleo.}
\end{questao}

\begin{questao}[2]{}
Um próton é abandonado em um campo uniforme de
$E=\SI{1e4}{\newton\per\coulomb}$. Determine sua aceleração.
\dados{$e=\pot{1.6}{-19}\,\si{\coulomb}$; $m_p=\pot{1.67}{-27}\,\si{\kilogram}$}
\resposta{$a=\pot{9.6}{11}\,\si{\meter\per\second\squared}$}
\resolucao{
\[
a=\frac{eE}{m_p}=\frac{(\pot{1.6}{-19})(\pot{1}{4})}{\pot{1.67}{-27}}
=\pot{9.58}{11}\,\si{\meter\per\second\squared}.
\]
\textbf{Cerca de $10^{11}$ vezes $g$} --- a gravidade é completamente
desprezível na dinâmica de partículas carregadas.\\
\textbf{Para o elétron}, a aceleração seria $1836$ vezes maior
($\pot{1.76}{15}\,\si{\meter\per\second\squared}$), pela razão das massas. É por
isso que, em tubos e aceleradores, são os elétrons que respondem primeiro.}
\end{questao}

\begin{questao}[2]{}
Duas cargas estão sobre o eixo $x$: $+\SI{4}{\micro\coulomb}$ em $x=0$ e
$-\SI{2}{\micro\coulomb}$ em $x=\SI{0.30}{\meter}$. Determine o campo em
$x=\SI{0.60}{\meter}$.
\resposta{$E=\SI{100}{\kilo\newton\per\coulomb}$, no sentido $-x$}
\resolucao{
\[
E_1=\frac{\pot{9}{9}(\pot{4}{-6})}{(0{,}60)^{2}}
=\pot{1}{5}\,\si{\newton\per\coulomb}\ (\text{sentido }+x);
\]
\[
E_2=\frac{\pot{9}{9}(\pot{2}{-6})}{(0{,}30)^{2}}
=\pot{2}{5}\,\si{\newton\per\coulomb}\ (\text{sentido }-x).
\]
\[
E=2\times10^{5}-1\times10^{5}=\pot{1}{5}\,\si{\newton\per\coulomb}\ (-x).
\]
\textbf{Observe:} a carga \emph{menor} produziu o campo \emph{maior}, por estar
duas vezes mais próxima --- e $1/r^{2}$ vence o fator $2$ de carga.
Compare os fatores: carga $\times\frac12$, distância$^{2}$
$\times\frac14$.}
\end{questao}

\begin{questao}[2]{}
Duas cargas iguais $Q=+\SI{2}{\micro\coulomb}$ estão em
$(\pm\SI{0.30}{\meter},\,0)$. Determine o campo no ponto
$(0,\ \SI{0.40}{\meter})$.
\resposta{$E=\SI{115.2}{\kilo\newton\per\coulomb}$, ao longo de $+y$}
\resolucao{Cada carga dista $r=\sqrt{0{,}30^{2}+0{,}40^{2}}=\SI{0.5}{\meter}$
(triângulo $3$--$4$--$5$):
\[
E_1=E_2=\frac{\pot{9}{9}(\pot{2}{-6})}{0{,}25}
=\pot{7.2}{4}\,\si{\newton\per\coulomb}.
\]
\textbf{Por simetria}, as componentes horizontais se cancelam e as verticais se
somam, cada uma multiplicada por $\cos\theta=0{,}40/0{,}50=0{,}8$:
\[
E=2(\pot{7.2}{4})(0{,}8)=\pot{1.152}{5}\,\si{\newton\per\coulomb}.
\]
\textbf{Uso da simetria:} reconhecê-la elimina metade do trabalho. Sempre que a
configuração for simétrica em relação ao ponto de interesse, identifique quais
componentes se anulam \emph{antes} de calcular.}
\end{questao}

\begin{questao}[2]{}
Duas placas paralelas separadas por $\SI{2}{\milli\meter}$ estão submetidas a
$\SI{600}{\volt}$. Determine o campo entre elas.
\resposta{$E=\SI{300}{\kilo\volt\per\meter}$}
\resolucao{
\[
E=\frac{U}{d}=\frac{600}{\pot{2}{-3}}=\pot{3}{5}\,\si{\volt\per\meter}.
\]
\textbf{Campo uniforme:} entre placas extensas e próximas, $E$ é praticamente o
mesmo em todos os pontos --- e independe da posição. É a configuração de
referência para acelerar e defletir partículas.\\
\textbf{Verificação de segurança:} $\SI{300}{\kilo\volt\per\meter}$ é um décimo
da rigidez do ar ($\SI{3}{\mega\volt\per\meter}$) --- há folga confortável.
Reduzir a separação para $\SI{0.2}{\milli\meter}$ com a mesma tensão provocaria
ruptura.}
\end{questao}

\begin{questao}[2]{}
Uma placa condutora extensa tem densidade superficial de carga
$\sigma=\SI{2}{\micro\coulomb\per\meter\squared}$. Determine o campo próximo à
sua superfície.
\dados{$\varepsilon_0=\pot{8.85}{-12}\,\si{\farad\per\meter}$}
\resposta{$E=\SI{226}{\kilo\newton\per\coulomb}$}
\resolucao{Para um condutor em equilíbrio,
\[
E=\frac{\sigma}{\varepsilon_0}=\frac{\pot{2}{-6}}{\pot{8.85}{-12}}
=\pot{2.26}{5}\,\si{\newton\per\coulomb}.
\]
\textbf{Fato notável:} o campo \textbf{não depende da distância} à placa (desde
que ela seja extensa e o ponto próximo). Diferentemente da carga puntiforme, não
há queda com $1/r^{2}$.\\
\textbf{Cuidado com a fórmula:} para um \emph{plano isolante} com carga nas duas
faces, o campo vale $\sigma/2\varepsilon_0$ --- metade. A diferença é que no
condutor toda a carga fica em uma face, e o campo interno é nulo.}
\end{questao}

\begin{questao}[2]{}
Compare o campo produzido por $Q_1=+\SI{9}{\micro\coulomb}$ a
$\SI{30}{\centi\meter}$ com o de $Q_2=+\SI{4}{\micro\coulomb}$ a
$\SI{20}{\centi\meter}$.
\resposta{$E_1=E_2=\SI{900}{\kilo\newton\per\coulomb}$ --- iguais}
\resolucao{
\[
E_1=\frac{\pot{9}{9}(\pot{9}{-6})}{0{,}09}=\pot{9}{5};
\qquad
E_2=\frac{\pot{9}{9}(\pot{4}{-6})}{0{,}04}=\pot{9}{5}\,\si{\newton\per\coulomb}.
\]
\textbf{Coincidência não é acaso:} as razões $Q/r^{2}$ são iguais, pois
$9/9=4/4$. Sempre que $Q\propto r^{2}$, o campo se mantém.\\
\textbf{Consequência:} se essas duas cargas estivessem em lados opostos de um
ponto, o campo resultante ali seria \emph{nulo} --- e essa é exatamente a
condição do "ponto de campo nulo" entre cargas de mesmo sinal.}
\end{questao}

\begin{questao}[2]{}
Uma carga em um ponto sofre campos de $\SI{3}{\kilo\newton\per\coulomb}$ na
direção $x$ e $\SI{4}{\kilo\newton\per\coulomb}$ na direção $y$, produzidos por
duas fontes distintas. Determine o campo resultante.
\resposta{$E=\SI{5}{\kilo\newton\per\coulomb}$, a $\ang{53.1}$ do eixo $x$}
\resolucao{Sendo perpendiculares, compõem-se por Pitágoras:
\[
E=\sqrt{3^{2}+4^{2}}=\SI{5}{\kilo\newton\per\coulomb};
\qquad
\theta=\arctan\frac{4}{3}=\ang{53.1}.
\]
\textbf{Contraste com o potencial:} se esses mesmos pontos tivessem potenciais
de $3$ e $\SI{4}{\kilo\volt}$, o total seria simplesmente
$\SI{7}{\kilo\volt}$ --- soma algébrica, sem geometria.\\
É essa diferença que torna o potencial a ferramenta preferida para cálculos, e o
campo a ferramenta preferida para entender forças e movimento.}
\end{questao}

\begin{questao}[2]{}
Em uma representação por linhas de campo, o que indica uma região onde as
linhas estão mais próximas entre si?
\begin{alternativas}[2]
\item potencial maior
\item campo mais intenso
\item carga maior naquele ponto
\item ausência de campo
\end{alternativas}
\resposta{(b)}
\resolucao{A \textbf{densidade de linhas} é proporcional ao módulo do campo ---
essa é a convenção que dá utilidade quantitativa ao desenho.\\
\textbf{Três regras das linhas de campo:}
\begin{enumerate}
\item Nascem em cargas positivas e morrem em negativas (ou no infinito).
\item São tangentes a $\vec E$ em cada ponto.
\item \textbf{Nunca se cruzam} --- em um cruzamento, o campo teria duas
direções, o que é impossível.
\end{enumerate}
\textbf{Cuidado com (a):} potencial e campo são grandezas distintas. Linhas
próximas indicam campo intenso, o que corresponde a equipotenciais próximas ---
mas não diz nada sobre o \emph{valor} de $V$.}
\end{questao}

\begin{questao}[2]{}
Uma carga produz campo de $\SI{8}{\kilo\newton\per\coulomb}$ no vácuo. O
conjunto é imerso em um dielétrico de $\varepsilon_r=4$. Determine o novo campo.
\resposta{$E=\SI{2}{\kilo\newton\per\coulomb}$}
\resolucao{
\[
E_{meio}=\frac{E_{vácuo}}{\varepsilon_r}=\frac{8}{4}
=\SI{2}{\kilo\newton\per\coulomb}.
\]
\textbf{Mecanismo:} o dielétrico se \emph{polariza} --- suas moléculas se
orientam e criam um campo interno oposto, que cancela parcialmente o campo
original.\\
\textbf{Consequência tecnológica:} o dielétrico reduz o campo para a mesma
carga, o que permite armazenar \emph{mais} carga para a mesma tensão --- é
exatamente por isso que $C=\varepsilon_r C_0$ em capacitores.}
\end{questao}

\begin{questao}[2]{}
Uma gota de massa $\SI{2}{\milli\gram}$ e carga $\SI{1}{\nano\coulomb}$ deve
ficar suspensa em repouso. Determine o campo elétrico necessário.
\dados{$g=\SI{9.8}{\meter\per\second\squared}$}
\resposta{$E=\SI{19.6}{\kilo\newton\per\coulomb}$}
\resolucao{Equilíbrio entre peso e força elétrica:
\[
qE=mg
\;\Longrightarrow\;
E=\frac{mg}{q}=\frac{(\pot{2}{-6})(9{,}8)}{\pot{1}{-9}}
=\pot{1.96}{4}\,\si{\newton\per\coulomb}.
\]
\textbf{Orientação:} o campo deve produzir força \emph{para cima}. Se a gota for
positiva, $\vec E$ aponta para cima; se negativa, para baixo.\\
\textbf{Este é o princípio do experimento de Millikan}, que determinou a carga
elementar. Medindo $E$ e $m$ de gotas suspensas, ele encontrou que todas as
cargas eram múltiplas inteiras de $\pot{1.6}{-19}\,\si{\coulomb}$ --- a
primeira evidência direta da quantização.}
\end{questao}

\blocoaprofundamento

\begin{questao}[3]{}
Duas cargas iguais $Q=+\SI{1}{\micro\coulomb}$ ocupam
$(\pm\SI{0.20}{\meter},\,0)$.
\begin{itensq}
\item Determine o campo em $(0,\ \SI{0.20}{\meter})$.
\item Determine o campo no ponto médio e justifique.
\end{itensq}
\resposta{(a) $E=\SI{159}{\kilo\newton\per\coulomb}$ ao longo de $+y$;
(b) $\vec E=\vec0$}
\resolucao{(a) Cada carga dista
$r=\sqrt{2}\,(0{,}20)=\SI{0.283}{\meter}$:
\[
E_1=E_2=\frac{\pot{9}{9}(\pot{1}{-6})}{0{,}08}
=\pot{1.125}{5}\,\si{\newton\per\coulomb}.
\]
As componentes em $x$ cancelam; as em $y$ somam-se, cada uma multiplicada por
$\cos\ang{45}=0{,}707$:
\[
E=2(\pot{1.125}{5})(0{,}707)=\pot{1.59}{5}\,\si{\newton\per\coulomb}.
\]
(b) \textbf{No ponto médio}, as duas contribuições têm módulos iguais e sentidos
\emph{opostos} (cada carga empurra para longe de si):
\[
\vec E=\vec 0.
\]
\textbf{Mas o potencial ali não é nulo:}
$V=2\ke Q/0{,}20=\SI{90}{\kilo\volt}$. É o caso complementar ao do dipolo, em
que $V=0$ e $\vec E\ne\vec0$ --- e mostra novamente que as duas condições são
independentes.\\
\textbf{Estabilidade:} uma carga positiva no ponto médio está em equilíbrio, mas
\emph{instável} na direção perpendicular ao eixo --- conforme o teorema de
Earnshaw.}
\end{questao}

\begin{questao}[3]{}
Um dipolo é formado por $+q$ e $-q$ separadas por $2a$. Considere um ponto do
eixo perpendicular à linha que as une, a distância $y$ do centro.
\begin{itensq}
\item Deduza a expressão exata do campo.
\item Obtenha a forma assintótica para $y\gg a$.
\item Compare a queda com a de uma carga isolada.
\end{itensq}
\resposta{$E=\dfrac{2\ke qa}{(y^{2}+a^{2})^{3/2}}\to\dfrac{2\ke qa}{y^{3}}$}
\resolucao{(a) Cada carga dista $r=\sqrt{y^{2}+a^{2}}$ e produz
$\ke q/r^{2}$. Por simetria, as componentes ao longo de $y$ se cancelam (uma
aponta para cima, a outra para baixo) e restam as componentes \emph{paralelas ao
eixo do dipolo}, cada uma multiplicada por $a/r$:
\[
E=2\cdot\frac{\ke q}{y^{2}+a^{2}}\cdot\frac{a}{\sqrt{y^{2}+a^{2}}}
=\frac{2\ke qa}{(y^{2}+a^{2})^{3/2}}.
\]
(b) Para $y\gg a$, $(y^{2}+a^{2})^{3/2}\to y^{3}$:
\[
E\approx\frac{2\ke qa}{y^{3}}=\frac{\ke p}{y^{3}},
\]
com $p=2qa$ o momento dipolar.\\
(c) \textbf{O campo do dipolo cai com $1/y^{3}$}, muito mais rápido que o
$1/r^{2}$ de uma carga isolada.\\
\textbf{Razão física:} de longe, as duas cargas quase se cancelam --- o sistema
é \emph{globalmente neutro}. O que sobra é o efeito de sua \emph{separação},
uma correção de ordem superior.\\
\textbf{Padrão geral:} carga isolada $\sim1/r^{2}$; dipolo
$\sim1/r^{3}$; quadrupolo $\sim1/r^{4}$. Quanto mais neutra a distribuição, mais
rápido o campo se extingue --- e é por isso que a matéria comum, neutra, não
exerce força elétrica a distância.}
\end{questao}

\begin{questao}[3]{}
Duas placas condutoras extensas e paralelas têm densidades superficiais
$+\sigma$ e $-\sigma$, com
$\sigma=\SI{1}{\micro\coulomb\per\meter\squared}$.
\begin{itensq}
\item Determine o campo entre as placas.
\item Determine o campo fora delas.
\item Justifique fisicamente.
\end{itensq}
\resposta{Entre: $E=\SI{113}{\kilo\newton\per\coulomb}$; fora: $E=0$}
\resolucao{(a) Cada plano contribui com $\sigma/2\varepsilon_0$, e
\textbf{entre} as placas as duas contribuições têm o \emph{mesmo} sentido (do
positivo para o negativo):
\[
E=2\cdot\frac{\sigma}{2\varepsilon_0}=\frac{\sigma}{\varepsilon_0}
=\frac{\pot{1}{-6}}{\pot{8.85}{-12}}=\pot{1.13}{5}\,\si{\newton\per\coulomb}.
\]
(b) \textbf{Fora}, as duas contribuições têm sentidos \emph{opostos} e módulos
iguais:
\[
E=\frac{\sigma}{2\varepsilon_0}-\frac{\sigma}{2\varepsilon_0}=0.
\]
(c) \textbf{Justificativa:} o campo de um plano infinito é uniforme e aponta
para longe dele (se positivo) ou para ele (se negativo). Na região externa,
qualquer ponto está do mesmo lado das duas placas --- e os campos se opõem.
Entre elas, o ponto está "depois" da positiva e "antes" da negativa, e os campos
se somam.\\
\textbf{Consequência: blindagem natural.} O par de placas confina completamente
seu campo na região interna --- não há campo escapando para fora. É por isso que
um capacitor de placas paralelas não interfere no circuito vizinho, e por que a
energia armazenada está inteiramente entre as armaduras.}
\end{questao}

\begin{questao}[3]{}
Uma esfera condutora de raio $R=\SI{10}{\centi\meter}$ tem carga
$Q=\SI{5}{\micro\coulomb}$.
\begin{itensq}
\item Determine o campo no interior, na superfície e a $\SI{30}{\centi\meter}$
      do centro.
\item Avalie se a configuração é fisicamente sustentável no ar.
\end{itensq}
\resposta{$0$; $\SI{4.5}{\mega\volt\per\meter}$;
$\SI{500}{\kilo\volt\per\meter}$ --- \textbf{insustentável}}
\resolucao{(a) \textbf{Interior:} $E=0$, condição de equilíbrio de qualquer
condutor.\\
\textbf{Superfície:} toda a carga se comporta como concentrada no centro, para
pontos externos:
\[
E=\frac{\ke Q}{R^{2}}=\frac{\pot{9}{9}(\pot{5}{-6})}{0{,}01}
=\pot{4.5}{6}\,\si{\volt\per\meter}.
\]
\textbf{A $\SI{30}{\centi\meter}$:}
$E=\dfrac{\pot{4.5}{4}}{0{,}09}=\pot{5}{5}\,\si{\volt\per\meter}$.\\
(b) \textbf{Insustentável.} O campo na superfície,
$\SI{4.5}{\mega\volt\per\meter}$, excede a rigidez dielétrica do ar
($\SI{3}{\mega\volt\per\meter}$) em $50\%$. O ar ioniza e a carga escapa por
descarga corona.\\
\textbf{Carga máxima real:}
$Q_{\max}=E_{rup}R^{2}/\ke=\SI{3.33}{\micro\coulomb}$. A esfera simplesmente
\emph{não consegue} reter os $\SI{5}{\micro\coulomb}$ enunciados.\\
\textbf{Lição de método:} sempre verifique se o resultado é fisicamente
possível. Um problema pode ser aritmeticamente consistente e descrever uma
situação inexistente.}
\end{questao}

\begin{questao}[3]{}
Um elétron parte do repouso em um campo uniforme de
$E=\SI{1000}{\volt\per\meter}$ e percorre $\SI{2}{\centi\meter}$.
\begin{itensq}
\item Determine a aceleração e o tempo de percurso.
\item Determine a velocidade final e a energia adquirida.
\end{itensq}
\dados{$e=\pot{1.6}{-19}\,\si{\coulomb}$; $m_e=\pot{9.11}{-31}\,\si{\kilogram}$}
\resposta{$a=\pot{1.76}{14}\,\si{\meter\per\second\squared}$;
$t=\SI{15.1}{\nano\second}$; $v=\pot{2.65}{6}\,\si{\meter\per\second}$;
$\SI{20}{\electronvolt}$}
\resolucao{(a)
\[
a=\frac{eE}{m_e}=\frac{(\pot{1.6}{-19})(1000)}{\pot{9.11}{-31}}
=\pot{1.756}{14}\,\si{\meter\per\second\squared};
\]
\[
t=\sqrt{\frac{2d}{a}}=\sqrt{\frac{2(0{,}02)}{\pot{1.756}{14}}}
=\pot{1.51}{-8}\,\si{\second}=\SI{15.1}{\nano\second}.
\]
(b) $v=at=(\pot{1.756}{14})(\pot{1.51}{-8})
=\pot{2.65}{6}\,\si{\meter\per\second}$.\\
\textbf{Energia:} o mais rápido é usar a tensão percorrida,
$U=Ed=1000(0{,}02)=\SI{20}{\volt}$:
\[
W=eU=\SI{20}{\electronvolt}=\pot{3.2}{-18}\,\si{\joule}.
\]
\emph{(Conferindo: $\frac12m_ev^{2}=\frac12(\pot{9.11}{-31})(\pot{2.65}{6})^{2}
=\pot{3.2}{-18}\,\si{\joule}$ \checkmark.)}\\
\textbf{Verificação relativística:} $v/c=0{,}9\%$ --- plenamente no regime
clássico.}
\end{questao}

\begin{questao}[3]{}
Compare a queda do campo com a distância para três configurações: carga
puntiforme, fio infinito carregado e plano infinito carregado.
\begin{itensq}
\item Escreva a dependência de cada uma.
\item Explique a origem geométrica das diferenças.
\end{itensq}
\resposta{$1/r^{2}$, $1/r$ e constante}
\resolucao{(a)
\begin{center}
\begin{tabular}{lcc}
\hline
Distribuição & Campo & Queda\\
\hline
Carga puntiforme & $\dfrac{\ke Q}{r^{2}}$ & $1/r^{2}$\\[6pt]
Fio infinito ($\lambda$) & $\dfrac{\lambda}{2\pi\varepsilon_0 r}$ & $1/r$\\[6pt]
Plano infinito ($\sigma$) & $\dfrac{\sigma}{2\varepsilon_0}$ & constante\\
\hline
\end{tabular}
\end{center}
(b) \textbf{Origem geométrica.} Pela lei de Gauss, o fluxo total é fixo pela
carga envolvida, e o campo é o fluxo dividido pela área da superfície:
\begin{itemize}
\item \textbf{Ponto:} a superfície é uma \emph{esfera}, de área
$4\pi r^{2}$ --- o campo cai com $1/r^{2}$.
\item \textbf{Fio:} a superfície é um \emph{cilindro}, de área lateral
$2\pi r\ell$ --- cresce apenas com $r$, e o campo cai com $1/r$.
\item \textbf{Plano:} a superfície é um \emph{prisma} cuja área \emph{não
depende} de $r$ --- e o campo não cai.
\end{itemize}
\textbf{Regra unificadora:} o campo cai como $1/r^{n}$, onde $n$ é o número de
dimensões em que a fonte é \emph{limitada}. Ponto: limitado em três
($n=2$). Fio: em duas ($n=1$). Plano: em uma ($n=0$).\\
\textbf{Consequência prática:} perto de uma placa carregada extensa, afastar-se
não ajuda --- o campo permanece. É o oposto da intuição formada com cargas
puntiformes.}
\end{questao}

\begin{questao}[3]{}
Uma haste isolante fina é carregada de modo que produz, em certo ponto, campo
de $\SI{600}{\newton\per\coulomb}$. A haste é então partida ao meio, e apenas
uma das metades permanece.
\begin{itensq}
\item O campo cai à metade? Justifique.
\item Determine em que condição isso valeria.
\end{itensq}
\resposta{Não necessariamente --- só vale se a metade removida contribuísse com
exatamente metade do campo}
\resolucao{(a) \textbf{Não, em geral.} O campo é a soma \emph{vetorial} das
contribuições de todos os elementos da haste, e cada elemento contribui com
módulo e direção diferentes, dependendo de sua distância e posição angular em
relação ao ponto.\\
Se o ponto estiver \emph{sobre a mediatriz} da haste, as duas metades
contribuem igualmente e o campo de fato cai à metade --- mas a \emph{direção}
resultante também muda, pois as componentes que antes se cancelavam agora não
se cancelam mais.\\
Se o ponto estiver \emph{sobre o prolongamento} da haste, a metade mais próxima
contribui muito mais que a distante, e remover a distante reduz o campo bem
menos que $50\%$.\\
(b) \textbf{A queda é exatamente à metade quando:}
\begin{itemize}
\item as duas metades contribuem com vetores \emph{idênticos} --- o que exige
simetria completa em relação ao ponto; \emph{e}
\item essa condição só se realiza em pontos muito distantes, onde a haste inteira
se comporta como carga puntiforme e a divisão é simplesmente $Q\to Q/2$.
\end{itemize}
\textbf{Lição:} a superposição vale \emph{sempre}, mas ela é vetorial. Reduzir
a fonte pela metade não reduz o campo pela metade, exceto em casos de simetria
particular.}
\end{questao}

\blocodesafio

\begin{questao}[4]{}
\textbf{(Campo nulo com potencial não nulo.)} Construa uma configuração em que
$\vec E=\vec0$ em um ponto onde $V\ne0$.
\begin{itensq}
\item Proponha o arranjo mais simples e determine os valores.
\item Explique por que isso é possível.
\item Determine se o inverso ($V=0$ com $\vec E\ne\vec0$) também ocorre.
\end{itensq}
\resposta{Duas cargas iguais $+Q$; no ponto médio $\vec E=\vec0$ e
$V=2\ke Q/a$}
\resolucao{(a) \textbf{Arranjo:} duas cargas iguais $Q=+\SI{2}{\micro\coulomb}$
em $(\pm\SI{0.20}{\meter},\,0)$. No ponto médio:
\[
\vec E=\vec 0
\quad\text{(contribuições opostas)};
\]
\[
V=2\cdot\frac{\ke Q}{0{,}20}
=2\cdot\frac{\pot{9}{9}(\pot{2}{-6})}{0{,}20}
=\pot{1.8}{5}\,\si{\volt}=\SI{180}{\kilo\volt}.
\]
Campo nulo, potencial de cento e oitenta mil volts.\\
(b) \textbf{Por que é possível.} As duas grandezas se relacionam por
$\vec E=-\nabla V$: o campo é a \emph{taxa de variação} de $V$, não seu valor.
Uma função pode ter valor grande e derivada nula --- é exatamente o que ocorre
em um ponto \emph{estacionário}.\\
No ponto médio, $V$ tem um mínimo ao longo do eixo que une as cargas: o valor é
alto, mas a inclinação é zero.\\
(c) \textbf{Sim, o inverso também ocorre.} Substituindo uma das cargas por
$-Q$, o ponto médio passa a ter $V=0$ (contribuições opostas se cancelam
algebricamente) e $\vec E\ne\vec0$ (os vetores apontam para o mesmo lado e se
somam).\\
\textbf{Síntese:} $V=0$ e $\vec E=\vec0$ são condições \textbf{logicamente
independentes}. Nenhuma implica a outra, em nenhum dos dois sentidos.}
\end{questao}

\begin{questao}[4]{}
\textbf{(Linhas de campo.)} Analise as linhas de campo de um dipolo elétrico.
\begin{itensq}
\item Descreva seu traçado.
\item Demonstre que linhas de campo nunca se cruzam.
\item Explique o que ocorre em pontos de campo nulo.
\end{itensq}
\resposta{Saem de $+q$ e entram em $-q$; não se cruzam porque $\vec E$ tem
direção única em cada ponto}
\resolucao{(a) \textbf{Traçado.} As linhas saem radialmente de $+q$, curvam-se e
entram em $-q$. A linha que segue o eixo do dipolo é reta; as demais formam
arcos progressivamente mais abertos. Algumas partem para o infinito e retornam.\\
A densidade de linhas é máxima entre as cargas --- onde o campo é mais intenso
--- e decresce com $1/r^{3}$ ao se afastar, conforme demonstrado antes.\\
(b) \textbf{Demonstração.} Suponha, por absurdo, que duas linhas se cruzem em
um ponto $P$. Por definição, a linha de campo é \emph{tangente} a $\vec E$ em
cada ponto. Se duas linhas distintas passam por $P$, elas têm tangentes
distintas ali --- e portanto $\vec E(P)$ teria \textbf{duas direções}
simultaneamente.\\
Mas $\vec E$ é uma função vetorial \emph{bem definida}: a cada ponto do espaço
corresponde \emph{um} vetor. Contradição. Logo, as linhas não se cruzam.\\
\textbf{Corolário:} o mesmo argumento vale para linhas de corrente, linhas de
campo magnético e qualquer outro campo vetorial.\\
(c) \textbf{Pontos de campo nulo são a exceção.} Onde $\vec E=\vec0$, não há
direção definida --- e ali as linhas \emph{podem} se encontrar, sem contradição.
São os \textbf{pontos de sela} do potencial.\\
Em um dipolo ($+q$ e $-q$) não existe tal ponto no espaço finito. Mas em duas
cargas \emph{iguais}, o ponto médio é exatamente um ponto de campo nulo --- e
ali as linhas se aproximam de todos os lados sem se tocarem, definindo uma
separatriz.}
\end{questao}

\begin{questao}[4]{}
\textbf{(Projeto de acelerador de placas.)} Projete um arranjo de placas
paralelas que imprima aceleração de
$a=\pot{2}{12}\,\si{\meter\per\second\squared}$ a um elétron.
\begin{itensq}
\item Determine o campo necessário.
\item Determine a tensão para separação de $\SI{1}{\centi\meter}$.
\item Determine o tempo de trânsito e a velocidade de saída.
\end{itensq}
\dados{$e=\pot{1.6}{-19}\,\si{\coulomb}$; $m_e=\pot{9.11}{-31}\,\si{\kilogram}$}
\resposta{$E=\SI{11.4}{\volt\per\meter}$; $U=\SI{0.114}{\volt}$;
$t=\SI{100}{\nano\second}$; $v=\pot{2}{5}\,\si{\meter\per\second}$}
\resolucao{(a)
\[
E=\frac{m_ea}{e}=\frac{(\pot{9.11}{-31})(\pot{2}{12})}{\pot{1.6}{-19}}
=\SI{11.4}{\volt\per\meter}.
\]
(b) $U=Ed=11{,}4(0{,}01)=\SI{0.114}{\volt}$.\\
(c)
\[
t=\sqrt{\frac{2d}{a}}=\sqrt{\frac{0{,}02}{\pot{2}{12}}}
=\pot{1}{-7}\,\si{\second}=\SI{100}{\nano\second};
\]
\[
v=at=(\pot{2}{12})(\pot{1}{-7})=\pot{2}{5}\,\si{\meter\per\second}.
\]
\textbf{O resultado é surpreendente e instrutivo.} Uma aceleração de
$\pot{2}{12}\,\si{\meter\per\second\squared}$ --- duzentos bilhões de vezes
$g$ --- exige apenas $\SI{0.114}{\volt}$, tensão menor que a de uma pilha
descarregada.\\
\textbf{Razão:} a razão carga/massa do elétron é enorme
($\pot{1.76}{11}\,\si{\coulomb\per\kilogram}$). Campos elétricos modestos
produzem acelerações colossais em partículas leves.\\
\textbf{Ressalva prática:} com tensão tão baixa, o projeto seria dominado por
efeitos parasitas --- diferenças de função-trabalho entre os metais das placas
(décimos de volt), cargas superficiais e ruído térmico. Um projeto real usaria
tensões de dezenas de volts e aceleração proporcionalmente maior.}
\end{questao}

\begin{questao}[4]{}
\textbf{(Cavidade em condutor.)} Uma carga puntiforme $+q$ é colocada no
interior de uma cavidade vazia dentro de um condutor \emph{neutro} e isolado.
\begin{itensq}
\item Determine a carga induzida em cada superfície.
\item Determine o campo dentro do metal e fora do condutor.
\item Analise o que muda se o condutor for aterrado.
\end{itensq}
\resposta{$-q$ na parede da cavidade, $+q$ na superfície externa; campo nulo no
metal e não nulo fora}
\resolucao{(a) \textbf{Parede da cavidade:} aplicando a lei de Gauss a uma
superfície inteiramente dentro do metal (onde $\vec E=\vec0$), o fluxo é nulo e
a carga envolvida também. Como a superfície envolve $+q$ e a parede interna:
\[
q+q_{interna}=0
\;\Longrightarrow\;
q_{interna}=-q.
\]
\textbf{Superfície externa:} o condutor era neutro, então a carga total deve
permanecer nula:
\[
q_{externa}=+q.
\]
(b) \textbf{Dentro do metal:} $\vec E=\vec0$, por equilíbrio.\\
\textbf{Fora:} a carga $+q$ da superfície externa produz campo --- e ele é
\textbf{esfericamente simétrico} se a superfície externa for esférica,
\emph{independentemente} de onde a carga esteja dentro da cavidade.\\
\textbf{Consequência notável:} de fora, é impossível saber a posição da carga
interna, nem a forma da cavidade. O condutor "apaga" essa informação.\\
(c) \textbf{Se o condutor for aterrado:} a carga $+q$ da superfície externa
escoa para a terra. Resta apenas $-q$ na parede da cavidade, e
\[
\vec E=\vec 0\quad\text{em todo o exterior}.
\]
\textbf{A blindagem torna-se completa nos dois sentidos.} É exatamente esse o
princípio da gaiola de Faraday aterrada --- e a diferença entre blindar contra
campos externos (que dispensa aterramento) e impedir que fontes internas
irradiem (que o exige).}
\end{questao}

\begin{questao}[4]{}
\textbf{(Esfera isolante uniformemente carregada.)} Uma esfera isolante maciça
de raio $R$ tem carga $Q$ distribuída uniformemente em seu \emph{volume}.
\begin{itensq}
\item Determine o campo para $r<R$ e para $r>R$.
\item Identifique onde o campo é máximo.
\item Compare com a esfera \emph{condutora} de mesma carga.
\end{itensq}
\resposta{$E=\dfrac{\ke Qr}{R^{3}}$ dentro e $\dfrac{\ke Q}{r^{2}}$ fora;
máximo na superfície}
\resolucao{(a) \textbf{Fora ($r>R$):} pela lei de Gauss com superfície esférica,
toda a carga está envolvida:
\[
E=\frac{\ke Q}{r^{2}},
\]
idêntico ao de uma carga puntiforme no centro.\\
\textbf{Dentro ($r<R$):} apenas a fração de carga interna à superfície gaussiana
conta. Como a densidade é uniforme, essa fração é $(r/R)^{3}$:
\[
E=\frac{\ke\,Q(r/R)^{3}}{r^{2}}=\frac{\ke Q\,r}{R^{3}},
\]
\textbf{proporcional a $r$} --- cresce linearmente do centro à superfície.\\
(b) \textbf{Máximo na superfície} ($r=R$), onde as duas expressões coincidem em
$E=\ke Q/R^{2}$. O campo é contínuo ali, embora com derivadas diferentes de
cada lado.\\
\emph{No centro, $E=0$ por simetria.}\\
(c) \textbf{Comparação com a esfera condutora.}
\begin{center}
\begin{tabular}{lcc}
\hline
Região & Isolante & Condutora\\
\hline
$r<R$ & $\ke Qr/R^{3}$ & $0$\\
$r=R$ & $\ke Q/R^{2}$ & $\ke Q/R^{2}$\\
$r>R$ & $\ke Q/r^{2}$ & $\ke Q/r^{2}$\\
\hline
\end{tabular}
\end{center}
\textbf{São indistinguíveis do lado de fora} --- a lei de Gauss só "enxerga" a
carga total envolvida, não sua distribuição. A diferença está inteiramente no
interior: no condutor a carga migra para a superfície e o campo interno é nulo;
no isolante ela permanece onde foi colocada.}
\end{questao}

\begin{questao}[4]{}
\textbf{(Traçado numérico de linhas de campo.)} Descreva um procedimento
computacional para traçar linhas de campo a partir de uma expressão conhecida de
$\vec E(x,y)$.
\begin{itensq}
\item Formule o algoritmo.
\item Indique como escolher os pontos de partida.
\item Aponte as dificuldades numéricas e como contorná-las.
\end{itensq}
\resposta{Integrar $\dfrac{\mathrm{d}\vec r}{\mathrm{d}s}
=\dfrac{\vec E}{|\vec E|}$ a partir de pontos distribuídos em torno das cargas}
\resolucao{(a) \textbf{Algoritmo.} A linha de campo é, por definição, tangente a
$\vec E$. Parametrizando pelo comprimento de arco $s$:
\[
\frac{\mathrm{d}\vec r}{\mathrm{d}s}=\frac{\vec E(\vec r)}{|\vec E(\vec r)|}.
\]
Numericamente, com passo $h$:
\begin{enumerate}
\item Escolha um ponto inicial $\vec r_0$.
\item Calcule $\vec E(\vec r_n)$ e normalize.
\item Avance: $\vec r_{n+1}=\vec r_n+h\,\hat E$ (Euler) ou, melhor, por
Runge--Kutta de quarta ordem.
\item Repita até sair do domínio ou atingir uma carga negativa.
\end{enumerate}
\textbf{Normalizar é essencial:} sem isso, o passo seria proporcional a
$|\vec E|$ e a linha "voaria" perto das cargas e travaria longe delas.\\
(b) \textbf{Pontos de partida.} Distribua-os em pequenos círculos ao redor de
cada carga \emph{positiva}, com espaçamento angular uniforme. Para que a
densidade de linhas represente corretamente a intensidade, o número de linhas
por carga deve ser \emph{proporcional a $|q|$} --- oito linhas para $+q$,
dezesseis para $+2q$.\\
(c) \textbf{Dificuldades e soluções.}
\begin{itemize}
\item \textbf{Singularidades:} $|\vec E|\to\infty$ nas cargas. Solução: iniciar
a uma distância pequena mas finita, e interromper a integração ao entrar em um
raio mínimo em torno de uma carga negativa.
\item \textbf{Pontos de campo nulo:} ali $\hat E$ é indefinido e o algoritmo
trava ou oscila. Solução: detectar $|\vec E|<\epsilon$ e encerrar a linha.
\item \textbf{Acúmulo de erro:} o método de Euler desvia sistematicamente em
regiões de curvatura alta. Solução: Runge--Kutta e passo adaptativo, menor onde
o campo varia rápido.
\item \textbf{Linhas que escapam:} defina um domínio e encerre ao ultrapassá-lo.
\end{itemize}
\textbf{Verificação do resultado:} as linhas traçadas devem ser perpendiculares
às equipotenciais calculadas independentemente --- teste que detecta quase todo
erro de implementação.}
\end{questao}

\begin{questao}[4]{}
\textbf{(Limite de campo em capacitor.)} Um capacitor de placas paralelas com
$d=\SI{0.1}{\milli\meter}$ usa dielétrico de rigidez
$\SI{20}{\kilo\volt\per\milli\meter}$ e $\varepsilon_r=3$.
\begin{itensq}
\item Determine a tensão máxima e o campo correspondente.
\item Determine a densidade superficial de carga no limite.
\item Discuta o compromisso entre capacitância e tensão.
\end{itensq}
\dados{$\varepsilon_0=\pot{8.85}{-12}\,\si{\farad\per\meter}$}
\resposta{$U_{\max}=\SI{2}{\kilo\volt}$;
$E=\SI{2e7}{\volt\per\meter}$;
$\sigma=\SI{531}{\micro\coulomb\per\meter\squared}$}
\resolucao{(a)
\[
U_{\max}=E_{rup}\,d=20\times0{,}1=\SI{2}{\kilo\volt};
\qquad
E=\pot{2}{7}\,\si{\volt\per\meter}.
\]
(b) No dielétrico, $E=\dfrac{\sigma}{\varepsilon_r\varepsilon_0}$, logo
\[
\sigma=\varepsilon_r\varepsilon_0E
=3(\pot{8.85}{-12})(\pot{2}{7})
=\pot{5.31}{-4}\,\si{\coulomb\per\meter\squared}.
\]
(c) \textbf{O compromisso.} A capacitância por unidade de área é
$C/A=\varepsilon_r\varepsilon_0/d$: reduzir $d$ \emph{aumenta} $C$. Mas
$U_{\max}=E_{rup}d$: reduzir $d$ \emph{diminui} a tensão suportável.\\
O produto revela o que realmente importa:
\[
\frac{C\,U_{\max}}{A}=\varepsilon_r\varepsilon_0E_{rup},
\]
que \textbf{não depende de $d$}. E a energia por unidade de volume,
\[
u=\frac12\varepsilon_r\varepsilon_0E_{rup}^{2},
\]
também não. \textbf{A geometria não decide nada} --- só o material.\\
\textbf{Consequência:} não existe truque de projeto que aumente a energia
armazenada por volume. Todo avanço em capacitores vem de materiais com
$\varepsilon_r$ ou $E_{rup}$ maiores --- e esses dois parâmetros costumam ser
antagônicos, o que torna o problema difícil.}
\end{questao}

\begin{questao}[4]{}
\textbf{(Superposição vetorial em anel.)} Um anel de raio $R$ tem carga
$Q$ uniformemente distribuída.
\begin{itensq}
\item Explique por que o campo é nulo no centro.
\item Determine, sem integrar, o campo em um ponto muito distante no eixo.
\item Discuta o que ocorre se metade do anel for removida.
\end{itensq}
\resposta{Centro: $\vec E=\vec0$ por simetria; longe: $E\to\ke Q/x^{2}$;
com meio anel, campo não nulo no centro}
\resolucao{(a) \textbf{No centro,} para cada elemento de carga existe outro
diametralmente oposto, à mesma distância $R$. Suas contribuições têm módulos
iguais e sentidos opostos, cancelando-se aos pares:
\[
\vec E_{centro}=\vec 0.
\]
\textbf{Note:} o \emph{potencial} no centro não é nulo --- vale
$V=\ke Q/R$, pois potenciais somam-se sem cancelamento entre contribuições de
mesmo sinal. Mais um caso de $\vec E=\vec0$ com $V\ne0$.\\
(b) \textbf{Muito longe} ($x\gg R$), todos os elementos estão praticamente à
mesma distância $x$ e praticamente na mesma direção. O anel se comporta como uma
carga puntiforme:
\[
E\to\frac{\ke Q}{x^{2}}.
\]
\textbf{Este é o teste que toda expressão deve passar:} qualquer distribuição
finita, vista de longe, deve reproduzir a lei de Coulomb.\\
(c) \textbf{Com meio anel}, a simetria de cancelamento se perde. Cada elemento
não tem mais seu par oposto, e o campo no centro \textbf{deixa de ser nulo}.\\
Pelo argumento da superposição ao contrário: como o anel inteiro dá campo nulo,
o campo do semianel restante é o \emph{negativo} do campo do semianel removido
--- ou seja, os dois têm o mesmo módulo. Integrando (como no caso do fio
semicircular), obtém-se
\[
E=\frac{2\ke Q'}{\pi R^{2}},
\]
com $Q'=Q/2$ a carga de cada metade, dirigido ao longo do eixo de simetria.}
\end{questao}
